\chapter{Conclusion and Future Work}

\section{Conclusion}
This project has successfully demonstrates an \textbf{algorithm-hardware codesign} for accelerating Vision Transformer (ViT) inference on FPGAs. Upon Addressing the core bottleneck—the: $O(N^2)$ cost of self-attention, we implemented an attention-based token pruning algorithm that dynamically removes low-importance tokens during inference.
Using Vitis HLS, we translated this idea into an efficient hardware architecture, achieving:
\begin{itemize}
    \item \textbf{Efficient Pruning Mechanism} using Attention scores
    \item \textbf{HLS-Friendly ``Skipping''} mechanisms for reduce computations
    \item \textbf{Performance Trade-off} for achieving better latency with minmal accuracy drop
\end{itemize}
Overall, this project provides a practical blueprint for deploying ViT models on resource-constrained FPGA platforms, showing that combining algorithmic optimization with custom hardware design is the key in overcoming high inference time for Vision Transformers.

\section{Future Work}
Future research directions upon this work could be:

\begin{itemize}
    % \item \textbf{Implementing a Hybrid Token Selector (HTS):} Our current work uses an Attention-based Token Selection (ATS) strategy. However, research shows that ATS can suffer from "immature or noisy attention scores" in the early layers. A key future enhancement would be to implement a \textbf{Keypoint-based Token Selection (KTS)} module. This mechanism would use a hardware-friendly algorithm like FAST keypoint detection on the raw image content to select tokens for the early layers. This would create a \textbf{Hybrid Token Selector (HTS)}, using KTS for the early layers and our existing ATS for the later layers, likely improving accuracy while still maintaining high GFLOPs savings.

   

    % \item \textbf{Introduce Convolutions (Hybrid ViT):} A primary weakness of ViTs is their lack of \textbf{"local inductive bias,"} which makes them data-hungry. Future work can be exploring a \textbf{Hybrid ViT architecture} by integrating \textbf{\textit{Convolutional operations}} into the transformer pipeline. These convolutions which can be added within or before the Transformer blocks enables the model to capture local spatial patterns in a much better way, where ViTs typically struggle with. This would also involve designing HLS-friendly "Convolutional Token Embedding" and "Convolutional Projection" modules, which can combine the local feature extraction of CNNs with the global context modeling of ViT Attention.
    \item \textbf{Deploying the Design on a Physical FPGA Board:}  
While this project presents a complete HLS-based architecture, running the design on an actual FPGA platform (such as the Xilinx Alveo or ZCU series boards) remains an important next step. A real hardware deployment would allow for precise measurement of on-board latency, resource utilization, and thermal behavior. Further optimizations such as custom memory hierarchies, pipelined attention blocks, and improved dataflow scheduling can then be explored to reduce latency even further.

\item \textbf{Adversarial attacks of Token Pruning:}  
Token pruning reduces computation by selectively removing less-important tokens—but it may also introduce new vulnerabilities. Future work could include a detailed study of \textbf{adversarial attacks on token pruning strategies}, exploring whether carefully perturbed inputs can mislead the pruning module into removing critical tokens. Studying these vulnerabilities and proposing adversarially robust pruning mechanisms (e.g., noise-resistant scoring functions, defense-aware pruning thresholds) would strengthen the reliability of ViT accelerators in security-sensitive applications.

 \item \textbf{Implement Token Merging (ToMe):} While the current project implements \textit{token pruning} which permanently removes a subset of tokens and can therefore lead to information loss, an alternative would be \textbf{Token Merging}. This maintains the original semantic content while still reducing the effective sequence length $N$ as computation progresses through the network. Since no token is completely removed, ToMe usually preserves accuracy significantly better than hard pruning.

\end{itemize}